% ============================
%   7. Future work
% ============================
\section{Remaining work and outlook}
\label{sec:future_work}

The progress documented in this report covers the first three milestones
of the planned development cycle. The remaining items, grouped by
priority, are listed below.

\subsection{Before the final thesis defense}

\begin{enumerate}
    \item \textbf{Flow-level features.} Aggregate packets into
        bidirectional 5-tuple flows and extend the feature vector with
        duration, packet count, byte count, mean inter-arrival, and
        protocol-specific fields. This is expected to improve the ML
        detector's behaviour on multi-packet attack patterns where a
        single packet is not informative on its own.
    \item \textbf{Evaluation on a public dataset.} Reproduce the
        training and evaluation pipeline on the
        CICIDS2017~\cite{sharafaldin2018cicids} dataset and on
        NSL-KDD~\cite{tavallaee2009detailed} to obtain comparable
        accuracy/precision/recall/F1 figures alongside the synthetic
        baseline already reported here.
    \item \textbf{Alembic migrations.} Replace
        \code{Base.metadata.create\_all()} with a versioned migration
        suite to support schema evolution without data loss.
    \item \textbf{Hardening of the live sniffer.} Document, and where
        possible automate, the Npcap installation on Windows and the
        capability flags on Linux (\code{cap\_net\_raw}).
    \item \textbf{Test coverage.} Expand the unit-test scaffolding
        (currently smoke tests on the signature engine and the
        dispatcher merge logic) towards an integration-test suite that
        replays a deterministic packet trace and asserts against the
        produced alerts.
\end{enumerate}

\subsection{Stretch goals}

\begin{enumerate}
    \item \textbf{Alert correlation.} Group consecutive alerts that
        share a source IP and category into a single ``incident'' view,
        analogous to the SecurityOnion ``case'' concept.
    \item \textbf{User-role separation.} Distinguish read-only analyst
        from admin in the UI; restrict rule editing to admins.
    \item \textbf{External rule import.} Translate a subset of Snort
        rules into the project's JSON match specification, providing a
        bridge to the larger rule ecosystem.
    \item \textbf{LLM-assisted triage.} A small chat panel on the
        Alert Detail page that summarises the captured packets in
        natural language --- a separate exploration tracked as a
        possible appendix chapter.
\end{enumerate}

\subsection{Risks and mitigations}

\begin{description}
    \item[Library compatibility on new Python releases.] Scientific
        wheels for Python~3.14 are not yet published on Windows. The
        launcher already refuses such interpreters and instructs the
        user to install Python~3.12. The risk is bounded.
    \item[Evaluation realism.] Synthetic datasets risk over-stating
        the ML detector's quality~\cite{sommer2010outside}. The planned
        CICIDS2017 evaluation is the mitigation.
    \item[Time budget.] The flow-level refactor is the largest single
        item remaining and is on the critical path. A reduced fallback
        (flow features computed only on a sliding window, no
        per-flow state in the database) is available if needed.
\end{description}

\subsection{Conclusion of this progress report}

The project has reached the point where the three engines --- signature,
machine-learning, and forensics --- cooperate inside a single,
self-hostable SIEM experience, accessible through a one-click launcher
suitable for academic demonstration. The remaining work is concentrated
on improving the ML detector's realism, evaluating it against
established public datasets, and tightening the production-readiness of
the deployment story. The codebase, the live demonstration of the
running application, and this document together constitute the
\textit{state of the art} of the bachelor's thesis at the date of this
report.
